
\documentclass{article}
\usepackage{colortbl}
\usepackage{makecell}
\usepackage{multirow}
\usepackage{supertabular}

\begin{document}

\newcounter{utterance}

\centering \large Interaction Transcript for game `ifeval', experiment `ifeval', episode 240 with refineonly{-}llama{-}3.1{-}8b{-}instruct.
\vspace{24pt}

{ \footnotesize  \setcounter{utterance}{1}
\setlength{\tabcolsep}{0pt}
\begin{supertabular}{c@{$\;$}|p{.15\linewidth}@{}p{.15\linewidth}p{.15\linewidth}p{.15\linewidth}p{.15\linewidth}p{.15\linewidth}}
    \# & \multicolumn{2}{c}{Player} && \multicolumn{2}{c}{Game Master} \\
    \hline

    \theutterance \stepcounter{utterance}  
    & & & \multicolumn{4}{p{0.6\linewidth}}{
        \cellcolor[rgb]{0.9,0.9,0.9}{
            \makecell[{{p{\linewidth}}}]{
                \texttt{\tiny{[P1$\langle$GM]}}
                \texttt{Write a casual blog post about how the outer solar system is different from the inner solar system, and what that means for the possibility of life. Wrap your entire response with double quotation marks. Your response should contain 17 or more sentences.} \\
            }
        }
    }
     \\ \\

    \theutterance \stepcounter{utterance}  
    & & \multicolumn{4}{p{0.6\linewidth}}{
        \cellcolor[rgb]{0.95,0.95,0.95}{
            \makecell[{{p{\linewidth}}}]{
                \texttt{\tiny{[P1$\leftrightarrow$P]}}
                \texttt{{[}DRAFT{]}} \\
\texttt{"Exploring the vastness of our solar system, we often find ourselves comparing the inner and outer regions. While the inner solar system is home to the rocky planets like Earth, Mercury, and Mars, the outer solar system is a realm of gas giants and icy worlds. The most notable difference between the two is the temperature and composition of the planets.} \\
\\ 
\texttt{In the inner solar system, the temperatures are relatively warm, with the average temperature on Earth being around 15°C (59°F). This warmth is due to the proximity to the sun, which provides the necessary energy for liquid water to exist. Liquid water is a crucial ingredient for life as we know it, making the inner solar system a prime candidate for hosting life.} \\
\\ 
\texttt{On the other hand, the outer solar system is a cold and dark place, with temperatures plummeting to {-}200°C ({-}330°F) or lower. The gas giants, such as Jupiter and Saturn, are massive balls of gas and liquid, with no solid surface to speak of. Their moons, like Europa and Enceladus, are icy worlds with subsurface oceans, but these oceans are thought to be frozen and inhospitable to life.} \\
\\ 
\texttt{However, the discovery of exoplanets in recent years has shown that the possibility of life existing in the outer solar system is not entirely ruled out. Some exoplanets have been found in the habitable zones of their respective stars, where temperatures are just right for liquid water to exist. These exoplanets are often referred to as "super{-}Earths" or "mini{-}Neptunes," and they offer a glimpse into what the outer solar system might have looked like in the distant past.} \\
\\ 
\texttt{The outer solar system is also home to a wealth of organic molecules, which are the building blocks of life. These molecules are thought to have originated from comets and asteroids that formed in the outer solar system and were later delivered to the inner solar system. The presence of these molecules on Earth has been linked to the origins of life, suggesting that the outer solar system played a crucial role in the emergence of life on our planet.} \\
\\ 
\texttt{In conclusion, while the inner solar system is more hospitable to life, the outer solar system is not entirely devoid of possibilities. The discovery of exoplanets and the presence of organic molecules suggest that life could exist in the outer solar system, albeit in forms that are very different from what we see on Earth. As we continue to explore the solar system, we may uncover more secrets about the origins of life and the possibility of life existing beyond our planet."} \\
            }
        }
    }
    & \\ \\

    \theutterance \stepcounter{utterance}  
    & \multicolumn{4}{p{0.6\linewidth}}{
        \cellcolor[rgb]{0.561,0.737,0.561}{
            \makecell[{{p{\linewidth}}}]{
                \texttt{\tiny{[P1$\rangle$GM]}}
                \texttt{"Exploring the vastness of our solar system, we often find ourselves comparing the inner and outer regions. While the inner solar system is home to the rocky planets like Earth, Mercury, and Mars, the outer solar system is a realm of gas giants and icy worlds. The most notable difference between the two is the temperature and composition of the planets, with the inner solar system boasting warmer temperatures and the outer solar system being a cold and dark place. The average temperature on Earth is around 15°C (59°F), while temperatures in the outer solar system plummet to {-}200°C ({-}330°F) or lower.} \\
\\ 
\texttt{This temperature difference is largely due to the proximity to the sun, which provides the necessary energy for liquid water to exist in the inner solar system. Liquid water is a crucial ingredient for life as we know it, making the inner solar system a prime candidate for hosting life. On the other hand, the gas giants, such as Jupiter and Saturn, are massive balls of gas and liquid, with no solid surface to speak of. Their moons, like Europa and Enceladus, are icy worlds with subsurface oceans, but these oceans are thought to be frozen and inhospitable to life.} \\
\\ 
\texttt{However, the discovery of exoplanets in recent years has shown that the possibility of life existing in the outer solar system is not entirely ruled out. Some exoplanets have been found in the habitable zones of their respective stars, where temperatures are just right for liquid water to exist. These exoplanets are often referred to as "super{-}Earths" or "mini{-}Neptunes," and they offer a glimpse into what the outer solar system might have looked like in the distant past.} \\
\\ 
\texttt{The outer solar system is also home to a wealth of organic molecules, which are the building blocks of life. These molecules are thought to have originated from comets and asteroids that formed in the outer solar system and were later delivered to the inner solar system. The presence of these molecules on Earth has been linked to the origins of life, suggesting that the outer solar system played a crucial role in the emergence of life on our planet.} \\
\\ 
\texttt{In addition, the outer solar system's unique environment may have given rise to life forms that are very different from what we see on Earth. For example, the high{-}energy radiation and extreme temperatures on some of the outer solar system's moons could have led to the evolution of life forms that are capable of surviving in these conditions. As we continue to explore the solar system, we may uncover more secrets about the origins of life and the possibility of life existing beyond our planet, and the discovery of exoplanets and organic molecules in the outer solar system is a crucial step in this journey."} \\
            }
        }
    }
    & & \\ \\

    \theutterance \stepcounter{utterance}  
    & & & \multicolumn{2}{p{0.3\linewidth}}{
        \cellcolor[rgb]{0.9,0.9,0.9}{
            \makecell[{{p{\linewidth}}}]{
                \texttt{\tiny{[GM$|$GM]}}
                {'startend:quotation': {}, 'length_constraints:number_sentences': {'relation': 'at least', 'num_sentences': 17}}
            }
        }
    }
    & & \\ \\

    \theutterance \stepcounter{utterance}  
    & & & \multicolumn{2}{p{0.3\linewidth}}{
        \cellcolor[rgb]{0.9,0.9,0.9}{
            \makecell[{{p{\linewidth}}}]{
                \texttt{\tiny{[GM$|$GM]}}
                \texttt{game\_result = WIN} \\
            }
        }
    }
    & & \\ \\

\end{supertabular}
}

\end{document}
